\section{Zadanie 9}
Niech \((X_n)^{\infty}_{n=1}\) będą niezależnymi zmiennymi losowymi z rozkładu wykładniczego z parametrem 1.
Dla danego \(k\) niech \(N(k) = min\{ n : \sum_{i=1}^{n} X_i > k\}\).
Wyznacz \(E(N(k))\) na dwa sposoby:
\begin{enumerate}[label=(\roman*)]
	\item\label{set09.task09.01} bezpośrednio korzystając z gęstości wyliczonych w jednym z poprzednich zadań
	      i wzoru \(E(N(k)) = \sum_{j=1}^{\infty} P(N(k) \geq j)\);
	\item\label{set09.task09.02} nie licząc ani jednej całki
	      i korzystając ze wzoru \(E(\sum_{i=1}^{N(k)}X_i) = E(N(k))\) (wcześniej go dowodząc).
\end{enumerate}

\subsection*{Rozwiązanie podpunktu \ref{set09.task09.01}}
Zauważmy, że \(N(k)\) to indeks pierwszego sumowania, które przekracza \(k\).
Niech \(S_n = \sum_{i=1}^{n} X_i\), wtedy
\[
	N(k) \geq j \Leftrightarrow S_{j-1} \leq k
\]
Z treści zadania 5, dla \(\lambda = 1\) mamy, że suma \(k\) niezależnych zmiennych o rozkładzie wykładniczym z parametrem \(\lambda\)
ma rozkład \(\Gamma(k, 1)\). \\
Więc gęstość zmiennej \(S_m\) (dla \(m \geq 1\)) wynosi(fakt 3, zestaw 9):
\[
	f_{S_m}(x) = \frac{1}{(m-1)!} x ^ {m-1} e^{-x}  \qquad \text{dla } x > 0
\]

Policzmy wartość oczekiwaną:
\[    E(N(k)) = \sum_{j=1}^{\infty} P(N(k) \geq j)\]
\[    E(N(k)) = \sum_{j=1}^{\infty} P(S_{j-1} \leq k)\]
\[    E(N(k)) = P(S_0 \leq k) + \sum_{j=2}^{\infty} P(S_{j-1} \leq k)\]
Pierwszy składnik jest równy 1 (bo \(S_0 = 0 \leq k\)). Dla \(j \geq 2\) podstawiamy \(m = j-1\)
\[    E(N(k)) = 1 + \sum_{m=1}^{\infty} \int_{0}^{k} \frac{x^{m-1}}{(m-1)!} e^{-x} dx \]
Teraz robimy sztuczkę i zamieniamy całkę z sumą:
\[    E(N(k)) = 1 + \int_{0}^{k} \sum_{m=1}^{\infty} \left( \frac{x^{m-1}}{(m-1)!} \right) e^{-x} dx \]
O to w nawiasie to szereg Taylora:
\[    E(N(k)) = 1 + \int_{0}^{k} e^x e^{-x} dx = 1 + \int_{0}^{k}dx = k + 1 \]

\subsection*{Rozwiązanie podpunktu \ref{set09.task09.02}}
\subsubsection*{Dowód tożsamości \(E(\sum_{i=1}^{N(k)}X_i) = E(N(k))\)}
Czyli korzystając ze zdefiniowanego \(S_n\) z podpunktu \ref{set09.task09.01} chcemy pokazać, że:
\[    E(S_{N(k)}) = E(N(k)) \]

Możemy zapisać sumę przy użyciu indykatorów:
\[ I_j = \begin{cases}
		1 \qquad N(k) \ge j \\
		0 \qquad \text{wpp}
	\end{cases}\]
\[ \sum_{i=1}^{N(k)} X_i = \sum_{j=1}^\infty X_j \cdot I_j \]
Korzystając z liniowości:
\[ E\left(\sum_{i=1}^{N(k)} X_i\right) = \sum_{j=1}^\infty E\left(X_j \cdot I_j\right) \]
Zauważmy, że zdarzenie \(N(k) \ge j\) oznacza, że proces nie zatrzymał się przed krokiem \(j\), czyli \(S_{j-1} \le k\).
Zdarzenie to zależy wyłącznie od zmiennych \(X_1, \dots, X_{j-1}\), a zatem jest niezależne od \(X_j\).
Możemy więc rozdzielić wartość oczekiwaną z prawej strony:
\[ E\left(X_j \cdot I_j\right) = E(X_j) \cdot P(N(k) \ge j) \]
Ponieważ \(X_j \sim \text{Exp}(1)\), to \(E(X_j) = 1\).
Więc:
\[ E\left(\sum_{i=1}^{N(k)} X_i\right) = 1 \cdot \sum_{j=1}^\infty P(N(k) \ge j) \]
Prawa strona to dokładnie \(E(N(k))\). Udowodniliśmy więc, że:
\[ E(S_{N(k)}) = E(N(k)) \]

\subsubsection*{Wyznaczenie wartości oczekiwanej}

Rozważmy lewą stronę równości, czyli \(E(S_{N(k)})\).
Zmienna \(S_{N(k)}\) oznacza wartość sumy w momencie pierwszego przekroczenia progu \(k\).
Możemy ją zapisać jako sumę progu \(k\) i nadwyżkę:
\[ S_{N(k)} = k + (S_{N(k)} - k) \]
Z własności braku pamięci rozkładu wykładniczego wynika, że wielkość,
o którą suma przekroczy poziom \(k\), ma ten sam rozkład co pojedynczy składnik sumy.
Zatem zmienna losowa \((S_{N(k)} - k)\) ma rozkład \(\text{Exp}(1)\), a wartość oczekiwana nadwyżki wynosi 1:
\[ E(S_{N(k)} - k) = 1 \]
Stąd otrzymujemy:
\[ E(S_{N(k)}) = k + 1 \]
Podstawiając ten wynik do udowodnionej tożsamości:
\[ E(N(k)) = E(S_{N(k)}) = k + 1  \]

\subsection*{Odpowiedź}
\[
	\boxed{E(N(k)) = k + 1}
\]
